% Generated draft: requires mathematical and source review.
\documentclass[11pt]{article}
\usepackage[margin=25mm]{geometry}
\usepackage{fontspec}
\setmainfont{Latin Modern Roman}
\usepackage{amsmath,amssymb,mathtools,unicode-math}
\setmathfont{Latin Modern Math}
\usepackage{enumitem,needspace,xurl,hyperref}
\setlength{\parindent}{0pt}
\setlength{\parskip}{5pt}
\setlength{\emergencystretch}{3em}
\newcommand{\N}{\mathbb{N}}
\newcommand{\Z}{\mathbb{Z}}
\newcommand{\Q}{\mathbb{Q}}
\newcommand{\R}{\mathbb{R}}
\newcommand{\C}{\mathbb{C}}
\newcommand{\F}{\mathbb{F}}
\newcommand{\A}{\mathbb{A}}
\newcommand{\PP}{\mathbb{P}}
\newcommand{\E}{\mathbb{E}}
\newcommand{\eps}{\varepsilon}
\newcommand{\dd}{\,\mathrm{d}}
\begin{document}
\section*{UnsolvedMath ID: 30006044}
\textbf{Source problem code:} catalogue entry
\subsection{Definitions and mathematical statement}
Let \(n\in\N\) with \(n\geq3\), and let \(\Lambda\in\R\) with \(\Lambda\geq0\). For an \(n\)-dimensional Riemannian manifold \((M,g)\) and a vector field \(X\) on \(M\), define its metric-dual one-form \(\alpha=X^\flat\) by \(\alpha(Y)=g(X,Y)\) for every vector field \(Y\). The closedness condition is
\[
\dd\alpha=0.
\]
Closedness must not be replaced by the stronger requirement that \(\alpha\) be globally exact.

\textbf{Classification problem.} For every such \(n\) and \(\Lambda\), classify all compact horizon data \((M,g,X)\) satisfying the near-horizon quasi-Einstein equation referred to as equation~(1) in the source, together with \(\dd X^\flat=0\). The supplied statement describes these data as equivalently the compact static near-horizon geometries.

\textbf{Unresolved formulation qualifications.} Equation~(1) is not reproduced in the supplied record. Consequently, its tensor equation, constants, sign conventions, and normalization of \(\Lambda\) cannot be specified from the supplied evidence. The record also does not specify the boundary convention, connectedness assumptions, or equivalence relation for the classification. The stated equivalence with static near-horizon geometries is supplied source information, not independently established here. These omissions prevent a fully self-contained mathematical formulation.

\textbf{Surrounding source context.} The original excerpt also discusses whether spherical topology of \(M\), together with equation~(1), forces the existence of a Killing vector field. This is contextual material, not an additional spherical-topology hypothesis on the classification problem above. The excerpt then attributes to Theorem~1, reference~[13], the assertion that any \(n\)-dimensional compact Riemannian manifold \((M,g)\) with a non-gradient vector field \(X\) satisfying equation~(1) admits a Killing vector field \(K\) such that
\[
\mathcal L_Kg=0,\qquad [K,X]=0.
\]
Here \(\mathcal L_Kg=0\) defines the Killing condition, and \([K,X]\) is the Lie bracket. In the usual global sense, non-gradient means that no globally defined function \(f\) satisfies \(X=\nabla_g f\). The supplied excerpt does not identify reference~[13], specify whether the asserted Killing field must be nonzero, or clarify the precise meaning of spherical topology. This contextual theorem is retained as an attributed source assertion and does not resolve the classification problem.
\subsection{Short English statement}
For every \(\Lambda\geq0\) and integer horizon dimension \(n\geq3\), classify the compact static near-horizon geometries, described in the supplied record by near-horizon quasi-Einstein data with closed metric-dual one-form. The defining equation and several classification conventions are missing, so this remains a qualified statement-only draft requiring mathematical review.
\subsection{Sources}
\begin{itemize}
\item Local UnsolvedMath registry, permanent ID \(30006044\), exact source code \texttt{catalogue entry}, title \emph{Classification of Static Near-Horizon Geometries}; supplied complete source record from the Oberwolfach Reports Open Problems collection. The proposed year is recorded as 2024; no proposer is identified.
\item Supplied source citation: \emph{Mathematical Aspects of General Relativity (2025)}, Oberwolfach Reports, pp.~2033--2112, DOI \texttt{10.4171/owr/2024/36}. \url{https://doi.org/10.4171/owr/2024/36}. The supplied literature list labels this same reference \emph{Oberwolfach Report 36 (2024)}. The differing year labels are preserved without reconciliation.
\item Supplied literature source: Lucietti and Kunduri, \emph{Near-horizon geometries}. \url{https://doi.org/10.12942/lrr-2013-8}.
\item The surrounding excerpt attributes its Killing-field theorem to Theorem~1, reference~[13]; the bibliography entry for [13] is not supplied.
\item Supplied literature assessment dated \texttt{2026-08-22}: status \emph{open}; the curation states that the report explicitly describes the classification for nonnegative cosmological constant in the stated dimensions as open. This is supplied dated curation, not independently verified here. The background expressly calls for independent verification of time-sensitive status and cited claims.
\item The registry labels the statement \texttt{corrected\_verified} and reports that the official report was checked, the question isolated from surrounding exposition, and notation repaired. These are supplied verification claims, not a record of independent source inspection or mathematical review in this draft.
\end{itemize}
\end{document}
